% Nouveau EGA V (cote 25): archive cote 25, site page 159 (PDF page 160).
% typescript. Draft compared: none.
% AI-transcribed by Claude Opus 5.5 (claude-opus-5-5, adaptive thinking, effort high), Message Batches API, batch ega25b (msgbatch_01A7n5YuXubUPNK1zsepi8RG), 2026-11-13; not yet checked by eye.
\scanpage{25/159}
\mspage{\hand{2}}

singulier quadratique  non dégénéré" sin on tient à une terminologie
voisine de la terminologie \xout{xé} recue en car.0 ; de même, on pourra
parler de "singularité quadratique non dégénérée" d'un préschéma loc
noeth , resp. "géommétriquement  quadratique non dégénérée" quand on est
loc de type fini sur un corps k.  Par contre, la terminologie
"zéro singulier ordinaire" ou mieux "zéro quadratique ordinaire",
et de même "singularité quadratique ordinaire","singularité géométrique-%
ment quadratique ordinaire",  devra s'entendre (conformément à l'usage)
comme correspondant à la notion de forme quadratique lisse. D'ailleurs,
il peut être indiqué de remplacer le terme "ordinaire" par "élémentaire",
et d'étendre cette terminologie aux singularités pas nécessairement
quadratiques, mais de multiplicité quelconque. Tu me diras ton impression
à cet égard.   -    Il semble bien d'autre part que le t exte SUP soit
formellement correct, en particulier dans\hand{/}prop.7  la notion qui s'introduit
vraiment est celle de zéro supersingulier.  Cela n'empêche qu'il y aurait
lieu, au moins en remarque, d'introduire aussi \xout{$V(\emptyset)$xxxxxxx$^{\text{sinxxél}}$} le sous-%
préschéma  $V(\emptyset)^{\text{ult sing}}$  de $V(\emptyset)^{\text{sup sing}}$, correspondant à la considéra-%
tion des zéros \unsure{ultra-singuliers|ultra-singulaers}, i.e. \xout{xxx} qui sont singuliers sans être
quadratiques élémentaires; il se décrit essentiellement par le même
procédé que $V(\emptyset)^{\text{sup sing}}$  , en prenant le préschéma des zéros du
"discriminant corrigé"  introduit plus bas, au lieu du discriminant ordi-%
naire.  Donc les seules corrections à apporter à SUP semblent d'ordre
terminologique (et elles doivent porter également, évidemment, sur \xout{xxxxx}\inserted[typed]{la}
\inserted[typed]{terminologie}\hand{,}
introduite \hand{$\checkmark$}dans SUP p.10\hand{/}\inserted{ci-dessus}).  Par contre,  IV 23.9.2. est faux tel quel
sauf si k est de car $\neq 2$ \uline{ou} les composantes irréductibles de X  de
dimension paire; dans le cas général, il faut supprimer "$f_n$ satisfait
aux conditions équivalentes de 8.8., en particulier \deleted{\illegible{1}}" ; le reste de la

\mspage{\hand{159}}
